% Hafta 6 — Bütünlük denetimi (self-hashing)
% Derleme: py -3.12 tools/latex/derle.py docs/week-6/assets/h06-03-self-hashing.tex
\documentclass[tikz,border=8pt]{standalone}
\usepackage{cen429-cizim}
\begin{document}
\begin{tikzpicture}
  \node[baslik] (b) at (0,0) {Bütünlük denetimi (self-hashing)};
  \node[altbaslik] at (b.south west) {uygulama kendi baytlarını doğrular};
  \node[kutu=36mm, anchor=north west] (n1) at (0,-2.4)
    {\kb{İkili dosya}\\/ korunan bölge};
  \node[koyukutu=36mm, right=9mm of n1] (n2)
    {\kb{HMAC-SHA-256}\\gizli anahtarla};
  \node[uyarikutu=36mm, right=9mm of n2] (n3)
    {\kb{altın değer ==}\\hesaplanan?};
  \draw[ok] (n1.east) -- (n2.west);
  \draw[ok] (n2.east) -- (n3.west);
  \node[iyikutu=40mm, below=26mm of n3, xshift=-32mm, anchor=north] (n4)
    {\kb{Bütünlük TAMAM}};
  \node[kotukutu=40mm, right=24mm of n4] (n5)
    {\kb{YAMA ALGILANDI}\\$\rightarrow$ tepki};
  \draw[iyiok] (n3.south) -- node[oketiket, above=3pt, xshift=-3mm] {evet} (n4.north);
  \draw[kotuok] (n3.south) -- node[oketiket, above=3pt, xshift=3mm] {hayır} (n5.north);
  \node[fit=(n1)(n2)(n3)(n4)(n5), inner sep=0pt] (grp) {};
  \node[kotudip, text width=155mm, anchor=north west] at ($(grp.south west)+(0,-8mm)$)
    {Tek başına yetmez: denetimin kendisi de yamalanabilir $\rightarrow$ örtüşen denetim gerekir.};
\end{tikzpicture}
\end{document}
